\documentclass[11pt,a4paper]{article}

\usepackage[utf8]{inputenc}
\usepackage[T1]{fontenc}
\usepackage{lmodern}
\usepackage{amsmath,amssymb,amsthm}
\usepackage{booktabs}
\usepackage{tabularx}
\usepackage{hyperref}
\usepackage[margin=2.5cm]{geometry}
\usepackage{natbib}
\usepackage{graphicx}
\usepackage{enumitem}
\usepackage{float}
\usepackage{tcolorbox}
\tcbuselibrary{skins,breakable}

\newtheorem{theorem}{Theorem}
\newtheorem{proposition}[theorem]{Proposition}
\newtheorem{corollary}[theorem]{Corollary}
\newtheorem{lemma}[theorem]{Lemma}
\newtheorem{definition}[theorem]{Definition}
\newtheorem{remark}[theorem]{Remark}
\newtheorem{example}[theorem]{Example}
\newtheorem{axiom}{Axiom}

\newcolumntype{Y}{>{\raggedright\arraybackslash}X}

\newtcolorbox{domainbox}[1]{
  colback=#1!5!white, colframe=#1!60!black,
  fonttitle=\bfseries, title={},
  enhanced
}

\hypersetup{
  colorlinks=true,
  linkcolor=blue!60!black,
  citecolor=green!40!black,
  urlcolor=blue!60!black,
  pdftitle={Pre-Logical States and the Birth of Information: How Reasoning Operates Before Truth Exists},
  pdfauthor={J. Arturo Ornelas Brand},
  pdfsubject={Pre-logical states, adjacent possible, free energy principle},
  pdfkeywords={pre-logical states, quaternionic logic, adjacent possible, free energy principle, information, collapse, Boolean lattice, substrate independence}
}

\title{Pre-Logical States and the Birth of Information:\\
How Reasoning Operates Before Truth Exists}

\author{J. Arturo Ornelas Brand\\
\small Independent Researcher, Aguascalientes, M\'exico\\
\small \texttt{arturoornelas62@gmail.com}}

\date{April 13, 2026}

\begin{document}

\maketitle

\begin{abstract}
Classical logic---Boolean, fuzzy, modal---assumes that truth values exist
before connectives operate on them. But many consequential decisions
occur \emph{before truth exists}: a firefighter dispatched to an
unidentified highway obstruction, a lawyer confronting a case with no
precedent, a city planner choosing infrastructure whose consequences
will unfold over decades. We formalize these situations as
\textbf{pre-logical states} using the four axes of the quaternionic
operator framework. A pre-logical state is a point in
$V = [0,1] \times [-1,1]^3$ where the truth axis~$r$ is undefined---not
indeterminate (a value exists but is unknown) but genuinely indefinite
(no value has crystallized). We prove that the space of possibilities
before collapse is a Boolean lattice with exact cardinality, expansion
rate, and metric, satisfying all four defining properties of Kauffman's
adjacent possible. We establish that the transition from pre-logical to
logical states---\textbf{collapse}---is governed by a damped harmonic
oscillator that is exactly the free energy minimization of Friston's
active inference framework, with solenoidal coupling ($\sigma > 0$)
necessary for oscillatory behavior. Two independent routes to
oscillation---dynamical (solenoidal flow) and geometric (cross-coupled
free energy)---are derived. We show that the duality pre-logical/\allowbreak logical
generates \textbf{information} as its compound concept, connecting the
quaternionic framework to Shannon entropy as a special case.
Substrate independence is demonstrated across four domains: quantum
measurement, legal reasoning, emergency response, and urban planning.
\end{abstract}

\noindent\textbf{Keywords:} pre-logical states; quaternionic logic;
adjacent possible; free energy principle; information; collapse;
Boolean lattice; substrate independence

%======================================================================
\section{Introduction}
\label{sec:intro}
%======================================================================

Logic begins with truth values. Boolean logic assigns $\{0,1\}$; fuzzy
logic assigns $[0,1]$; modal logic quantifies over possible worlds where
truth is defined in each. Every formal system in the logical tradition
takes truth as \emph{given} and asks what follows from it.

But consider:
\begin{itemize}[nosep]
\item A firefighter receives a report: ``large object blocking Highway~45.''
  The object could be a fallen tree, an overturned vehicle, a dead animal,
  or a boulder. No truth value for ``the object is a tree'' exists---not
  because the firefighter doesn't know, but because the situation has not
  yet been observed. What equipment should the truck carry?
\item A lawyer faces a case with no precedent in the jurisdiction.
  The applicable legal rule is not merely unknown---it does not yet exist.
  It will be \emph{created} by the judge's ruling. How should the lawyer
  argue?
\item A city planner must allocate budget between water infrastructure and
  transit for the year 2040. The city's needs in 2040 are not uncertain in
  the probabilistic sense---they depend on decisions not yet made, by people
  not yet born, under conditions not yet imaginable.
\end{itemize}

In each case, the agent must reason \emph{before truth exists}. This is
not probabilistic uncertainty (where a truth value exists but is hidden
behind a distribution). It is a qualitatively different state: the truth
axis is \textbf{indefinite}---no value has crystallized.

The quaternionic operator framework \citep{p11} provides the formal
apparatus to distinguish these states. The operator
$O = \{0,1,+,i,j,k\}$ acts on a space $V = [0,1] \times [-1,1]^3$
with four axes: truth~($r$), potentiality~($i$), verification~($j$),
and selection~($k$). The G-lattice of connectives is defined on this
space, with the commutation theorem establishing that connectives commute
with the group action if and only if they depend solely on~$r$. The
companion paper on duality synthesis \citep{p12} shows how opposing
poles generate new truth through the synthesis operation
$a \oplus \bar{a} = \mathrm{lcm}(\Phi(a), \Phi(\bar{a}))$.

This paper addresses what happens \emph{before} the G-lattice can
operate---in the regime where $r$ is not yet assigned. We formalize
pre-logical states (\S\ref{sec:prelogical}), establish the Boolean
lattice structure of the possibility space (\S\ref{sec:lattice}), define
the operations valid before collapse (\S\ref{sec:operations}), derive
collapse as an algebraic operation (\S\ref{sec:collapse}), prove the
dynamical equivalence between collapse and free energy minimization
(\S\ref{sec:wave-fep}), identify two independent routes to oscillatory
collapse (\S\ref{sec:two-routes}), show that information emerges as
the compound concept of the pre-logical/\allowbreak logical duality
(\S\ref{sec:information}), and demonstrate substrate independence across
four domains (\S\ref{sec:substrate}).

\paragraph{Contributions.}
\begin{enumerate}[nosep]
\item Formal definition of pre-logical states (Definition~\ref{def:prelogical}).
\item Proof that the possibility space is a Boolean lattice satisfying all
  four properties of the adjacent possible
  (Theorems~\ref{thm:lattice}--\ref{thm:nonprestate}).
\item Four operations valid before truth exists
  (Propositions~\ref{prop:expansion-op}--\ref{prop:selection-op}).
\item Collapse as an algebraic map $\mathbf{j}: \mathrm{Im}(V) \to V$
  (Definition~\ref{def:collapse}, Theorem~\ref{thm:collapse-properties}).
\item Exact equivalence between the damped wave and FEP dynamics
  (Theorem~\ref{thm:wave-fep-equiv}).
\item Necessity of solenoidal coupling for oscillation
  (Remark~\ref{rem:solenoidal}).
\item Two independent routes to oscillation
  (Proposition~\ref{prop:geometric}).
\item Information as compound concept of the pre-logical/\allowbreak logical duality
  (Proposition~\ref{thm:information}).
\item Substrate independence across four domains (Section~\ref{sec:substrate}).
\end{enumerate}

%======================================================================
\section{Preliminaries}
\label{sec:prelim}
%======================================================================

We recall the necessary elements from the quaternionic operator
framework \citep{p11} and the prime algebra \citep{p1}.

\paragraph{The quaternionic operator.}
The operator $O = \{0, 1, +, i, j, k\}$ acts on
$V = [0,1] \times [-1,1]^3$, where each point
$v = (r, i, j, k) \in V$ represents:
\begin{itemize}[nosep]
\item $r \in [0,1]$: truth (real axis),
\item $i \in [-1,1]$: potentiality (imaginary axis),
\item $j \in [-1,1]$: verification (empirical axis),
\item $k \in [-1,1]$: recursive selection (decision axis).
\end{itemize}
The four axes satisfy the quaternionic relation $ij = k$:
potentiality composed with verification yields selection.

\paragraph{The G-lattice.}
The connectives $\{\wedge, \vee, \neg, \to\}$ form a residuated lattice
on~$V$ with the $\mathrm{SU}(2)$ group action
$\alpha_q(v) = qvq^{-1}$ for unit quaternions $q \in S^3$. The
\emph{commutation theorem} \citep{p11} states that a connective $c$
satisfies $c \circ \alpha_q = \alpha_q \circ c$ for all $q$ if and only
if $c$ depends only on~$r$. Connectives that touch the imaginary axes
are context-sensitive.

\paragraph{Prime algebra.}
Each concept $C$ is encoded as a product of primes
$\Phi(C) = \prod p_i^{a_i}$ drawn from a set of $k$ irreducible
attributes \citep{p1}. The prime factor count
$\omega(\Phi(C))$ measures constructive complexity. The operations
$\gcd$ (shared structure), $\mathrm{lcm}$ (combined structure), and
$\bmod$ (containment test) are the algebraic primitives.

\paragraph{Synthesis.}
The duality synthesis operation \citep{p12} combines opposing poles:
$a \oplus \bar{a} = \mathrm{lcm}(\Phi(a), \Phi(\bar{a}))$, producing
a compound concept whose complexity strictly exceeds that of either pole
when both contribute at least one distinguishing prime.

\paragraph{Free energy principle.}
Under the Laplace approximation with generalized coordinates
$\tilde{\mu} = (\mu, \mu')$, the Friston dynamics
\citep{friston2010,buckley2017} take the form:
\begin{equation}
\label{eq:friston}
\dot{\tilde{\mu}} = D\tilde{\mu}
  - (\Gamma + Q)\,\nabla_{\tilde{\mu}} F
\end{equation}
where $D$ is the shift operator, $\Gamma$ the dissipation matrix, $Q$
the solenoidal (antisymmetric) matrix, and $F$ the variational free
energy.

%======================================================================
\section{The Pre-Logical State}
\label{sec:prelogical}
%======================================================================

\begin{definition}[Pre-logical state]
\label{def:prelogical}
A \textbf{pre-logical state} is a point in the extended space
$\hat{V} = \{?\} \cup [0,1] \times [-1,1]^3$ of the form:
\begin{equation}
v_{\mathrm{pre}} = (?, \; i, \; j, \; k)
  \quad\text{with } i > 0 \text{ and } j < 1
\end{equation}
where $? \notin [0,1]$ denotes that the truth axis is \textbf{undefined}
---not a missing value (as in incomplete information) but a genuinely
non-existent value (nothing has crystallized to assign truth to).
\end{definition}

The conditions $i > 0$ and $j < 1$ are not arbitrary:
\begin{itemize}[nosep]
\item $i > 0$: at least some potentiality remains open. If $i = 0$,
  there are no possibilities to explore---the state is already determined
  (the dogma pathology of \citealt{p11}).
\item $j < 1$: verification is incomplete. If $j = 1$, the state has
  been fully verified and truth is assigned---we are already in the
  logical regime.
\end{itemize}

\begin{remark}[Indefinite vs.\ indeterminate]
\label{rem:indefinite}
The distinction between \emph{indefinite} and \emph{indeterminate} is
critical. An indeterminate truth value exists but is epistemically
inaccessible (``the coin has landed but I haven't looked'').
Probability theory handles indeterminacy. An indefinite truth value
does not exist (``the coin is still spinning''). No probability
distribution can be assigned because the sample space itself is open.
The pre-logical state formalizes indefiniteness, not indeterminacy.
\end{remark}

\begin{remark}[Relation to the four pathologies]
\label{rem:pathologies}
The pathologies identified in \citet{p11}---dogma ($i = 0$), delirium
($r = 0$), pseudoscience ($j = 0$), and stagnation ($k = 0$)---are
collapses of individual axes within the logical regime. The pre-logical
state is qualitatively different: it is not a collapse of $r$ to zero
(that would be delirium, where truth is denied) but the
\emph{absence} of $r$ altogether. The pre-logical state precedes the
regime where pathologies are defined.
\end{remark}

%======================================================================
\section{The Boolean Lattice of Possibilities}
\label{sec:lattice}
%======================================================================

Before truth crystallizes, the agent navigates a space of
\emph{possibilities}---potential truths that might emerge upon collapse.
We show that this space has the structure of a Boolean lattice.

\begin{definition}[Prime product space]
\label{def:pps}
Let $P = \{p_1, p_2, \ldots, p_k\}$ be a set of $k$ distinct primes
representing irreducible attributes. The \textbf{prime product space} is:
\begin{equation}
\mathcal{L}(P) = \left\{ \prod_{i \in S} p_i
  : S \subseteq \{1, 2, \ldots, k\} \right\}
\end{equation}
with partial order given by divisibility: $A \leq B \iff A \mid B$.
\end{definition}

\begin{theorem}[Boolean lattice]
\label{thm:lattice}
$(\mathcal{L}(P), \mid)$ is a Boolean lattice isomorphic to
$(2^P, \subseteq)$:
\begin{itemize}[nosep]
\item \emph{Isomorphism}: $\varphi(S) = \prod_{i \in S} p_i$ is a
  bijection (by the Fundamental Theorem of Arithmetic).
\item \emph{Meet}: $\gcd(A, B) = \varphi(\varphi^{-1}(A) \cap \varphi^{-1}(B))$.
\item \emph{Join}: $\mathrm{lcm}(A, B) = \varphi(\varphi^{-1}(A) \cup \varphi^{-1}(B))$.
\item \emph{Complement}: $\overline{A} = (\prod P) / A$.
\item \emph{Bottom}: $1$ (empty product). \emph{Top}: $\prod P$ (all primes).
\item \emph{Atoms}: individual primes $p_1, \ldots, p_k$.
\item \emph{Rank function}: $r(A) = \omega(A)$ (number of distinct prime
  factors).
\end{itemize}
\end{theorem}

\begin{proof}
The map $\varphi: 2^P \to \mathcal{L}(P)$ defined by
$\varphi(S) = \prod_{i \in S} p_i$ is a bijection by the Fundamental
Theorem of Arithmetic. The lattice operations follow from
$\min(a_i, c_i)$ and $\max(a_i, c_i)$ on exponents, which correspond to
intersection and union of the index sets when exponents are binary.
\end{proof}

\begin{definition}[Adjacent possible in the lattice]
\label{def:ap}
Given a set of known concepts $K \subseteq \mathcal{L}(P)$, the
\textbf{adjacent possible} is:
\begin{equation}
AP(K) = \{ x \in \mathcal{L}(P) \setminus K :
  \exists\, y \in K \text{ s.t.\ } d_H(x, y) = 1 \}
\end{equation}
where $d_H$ is the Hamming distance between binary signatures (number of
primes by which $x$ and $y$ differ).
\end{definition}

We prove that this lattice satisfies all four defining properties of
Kauffman's adjacent possible \citep{kauffman2000,loreto2017}.

\subsection{Property 1: Expansion}

\begin{theorem}[Exact doubling]
\label{thm:expansion}
Adding one new prime $p_{k+1}$ to the repertoire exactly doubles the space:
\begin{equation}
|\mathcal{L}(P \cup \{p_{k+1}\})| = 2 \cdot |\mathcal{L}(P)|
\end{equation}
\end{theorem}

\begin{proof}
$|\mathcal{L}(P)| = |2^P| = 2^k$. Adding $p_{k+1}$ gives
$|\mathcal{L}(P \cup \{p_{k+1}\})| = 2^{k+1} = 2 \cdot 2^k$. The new
elements are exactly $\{A \cdot p_{k+1} : A \in \mathcal{L}(P)\}$---every
existing concept paired with the new prime. The old elements remain.
\end{proof}

\begin{corollary}
\label{cor:exponential}
After discovering $n$ new primes sequentially, the space grows by factor
$2^n$. Growth is exponential in the number of discoveries.
\end{corollary}

\subsection{Property 2: Path dependence}

\begin{theorem}[Exploration depends on history]
\label{thm:path}
The set of explored concepts at time~$t$ depends on the order in which
concepts are explored, not just on which primes are available.
\end{theorem}

\begin{proof}
Let $P = \{2, 3, 5\}$ and suppose an agent explores one adjacent concept
per step from $K_0 = \{1\}$.

\emph{Path~$\alpha$}: choose $p = 2$ first. At $t = 1$:
$K = \{1, 2\}$, $AP = \{3, 5, 6, 10\}$. Choose $6 = 2 \cdot 3$:
$K_\alpha = \{1, 2, 6\}$.

\emph{Path~$\beta$}: choose $p = 5$ first. At $t = 1$:
$K = \{1, 5\}$, $AP = \{2, 3, 10, 15\}$. Choose $15 = 3 \cdot 5$:
$K_\beta = \{1, 5, 15\}$.

Both paths have the same $P$ and the same total space
$\mathcal{L}(P) = \{1, 2, 3, 5, 6, 10, 15, 30\}$. But at $t = 2$,
$K_\alpha \cap K_\beta = \{1\}$: the explored sets are disjoint except
for the bottom element. Their adjacent possibles also differ.
\end{proof}

\subsection{Property 3: Faster-than-exploration}

\begin{theorem}[Vanishing exploration ratio]
\label{thm:faster}
For any agent exploring at rate~$r$ (concepts per unit time) with primes
discovered at rate~$q > 0$ (primes per unit time), the exploration ratio
$\rho(t) \to 0$ as $t \to \infty$.
\end{theorem}

\begin{proof}
At time~$t$:
\begin{equation}
\rho(t) = \frac{|K(t)|}{|\mathcal{L}(t)|}
  \leq \frac{r \cdot t}{2^{k_0 + qt}}
\end{equation}
Since $rt$ grows linearly and $2^{k_0 + qt}$ grows exponentially,
$\rho(t) \to 0$ by L'H\^opital's rule.
\end{proof}

\subsection{Property 4: Non-prestatability}

\begin{theorem}[Future concepts are unknowable]
\label{thm:nonprestate}
Given known primes~$P$ and known concepts $K \subseteq \mathcal{L}(P)$,
concepts involving an undiscovered prime $p_{k+2}$ cannot be enumerated
at the time $p_{k+1}$ is discovered.
\end{theorem}

\begin{proof}
The concept $p_{k+1} \cdot p_{k+2} \in \mathcal{L}(P \cup \{p_{k+1}, p_{k+2}\})$
but $p_{k+1} \cdot p_{k+2} \notin \mathcal{L}(P \cup \{p_{k+1}\})$
because $p_{k+2} \notin P \cup \{p_{k+1}\}$. The identity of $p_{k+2}$
is unknown at $p_{k+1}$'s discovery. Therefore the full future space
cannot be prestated from the current state.
\end{proof}

\begin{remark}[Conservative non-prestatability]
\label{rem:conservative}
In Kauffman's chemical setting, non-prestatability is stronger: even the
\emph{types of reactions} possible with new molecules may be unknown.
In the prime lattice, the operation is always multiplication---only the
identity of future primes is unknown. This makes the prime version a
conservative model: the weakest form that still satisfies the property.
This is an advantage for formalization, as it isolates non-prestatability
from operational novelty.
\end{remark}

\begin{remark}[Empirical support]
\label{rem:zipf}
Random walks on the Boolean lattice of the first 20 primes
($2^{20}$ nodes), using additive reinforcement with novelty
parameter~$\nu = 5.0$, produce a near-exact Zipf law:
$\alpha_Z = 0.995$, $R^2 = 0.966$ (500 independent walks of 10,000
steps each; data from the empirical validation of the prime
algebra~\citep{p4}).
\end{remark}

%======================================================================
\section{Operations Before Collapse}
\label{sec:operations}
%======================================================================

In a pre-logical state, the standard connectives ($\wedge, \vee, \neg,
\to$) are undefined because they require truth values. But four
operations remain well-defined:

\begin{proposition}[Expansion]
\label{prop:expansion-op}
The operation $\mathrm{expand}: \mathcal{L}(P) \to \mathcal{L}(P \cup \{p_{k+1}\})$
that adds a new prime to the repertoire is well-defined in a pre-logical
state. It increases $i$ (more possibilities open), does not require $r$,
and doubles the cardinality of the possibility space
(Theorem~\ref{thm:expansion}).
\end{proposition}

\begin{proof}
The expansion operation acts on the lattice structure, not on truth
values. It extends the prime set $P$ to $P' = P \cup \{p_{k+1}\}$,
creating $|\mathcal{L}(P)|$ new elements. No truth assignment is
consulted.
\end{proof}

\begin{proposition}[Partial evaluation]
\label{prop:partial-eval}
The operation $j_\epsilon: \hat{V} \to \hat{V}$ that increments the
verification axis by $\epsilon > 0$ without resolving $r$ is
well-defined. It represents gathering evidence that narrows the
possibility space without collapsing it:
\begin{equation}
j_\epsilon(?, i, j, k) = (?, \; i - \delta, \; j + \epsilon, \; k)
\end{equation}
where $\delta \geq 0$ is the reduction in potentiality induced by the
evidence (some possibilities are eliminated).
\end{proposition}

\begin{proof}
The operation modifies only the $i$ and $j$ coordinates, which are
defined in pre-logical states. The constraint $i > 0$ is maintained as
long as $\delta < i$; the constraint $j < 1$ is maintained as long as
$\epsilon < 1 - j$. Neither condition references~$r$.
\end{proof}

\begin{proposition}[Minimax preparation]
\label{prop:minimax}
Given candidate collapse outcomes
$T = \{v_1, \ldots, v_m\} \subset V$, the \textbf{minimax preparation}
is the point $v^* \in V$ that minimizes the maximum distance to all
possible outcomes:
\begin{equation}
v^* = \arg\min_{v \in V} \max_{v_\ell \in T}
  \, \|v - v_\ell\|
\end{equation}
where $\|\cdot\|$ is the Euclidean norm on
$V = [0,1]\times[-1,1]^3 \subset \mathbb{R}^4$.
This operation is well-defined in a pre-logical state because it
operates on the geometry of~$V$, not on a specific truth value.
\end{proposition}

\begin{proof}
The Euclidean distance is continuous on $V$, and $V$ is compact (closed
and bounded in $\mathbb{R}^4$). The function
$v \mapsto \max_\ell \|v - v_\ell\|$ is continuous as the maximum of
finitely many continuous functions. Existence of $v^*$ follows from the
extreme value theorem. The computation requires only the candidates $T$
(which the agent can enumerate from the possibility lattice), not the
actual truth value.
\end{proof}

\begin{proposition}[Selection]
\label{prop:selection-op}
The operation $k_\delta: \hat{V} \to \hat{V}$ that commits to a
direction reduces the effective possibility space by pruning branches of
the lattice:
\begin{equation}
k_\delta(?, i, j, k) = (?, \; i - \delta', \; j, \; k + \delta)
\end{equation}
where $\delta' \geq 0$ is the reduction in potentiality caused by the
commitment. Selection differs from partial evaluation: it reduces
possibilities by \emph{choice} (closing doors), not by \emph{evidence}
(ruling out options).
\end{proposition}

\begin{proof}
The quaternionic relation $ij = k$ ensures that selection is defined
whenever both potentiality ($i > 0$) and some verification ($j > 0$ or
$j = 0$ with prior evidence) are present. Axiom~K3 of the quaternionic
framework \citep{p11} states that $i > 0 \wedge j > 0$ enables but does
not determine $k$: selection requires an agent's commitment beyond what
the evidence provides.
\end{proof}

\begin{remark}[Ordering of operations]
The four operations have a natural partial order. Expansion ($i$ increases)
typically precedes partial evaluation ($j$ increases, $i$ decreases), which
typically precedes selection ($k$ increases, $i$ decreases). But this
ordering is not strict: an agent may select ($k$) before fully evaluating
($j$), accepting higher risk for faster action. The firefighter who brings
all equipment is performing minimax; the one who sends a drone first is
performing partial evaluation; the one who picks equipment based on
experience is selecting.
\end{remark}

\begin{remark}[Operations as structural types]
\label{rem:operations-types}
The four operations specify what \emph{type} of action is available in a
pre-logical state: expanding possibilities, gathering evidence, preparing
for all outcomes, or committing. The parameters $\delta$ (potentiality
reduction from evidence) and $\delta'$ (potentiality reduction from
commitment) are domain-specific: how much one drone observation narrows
the firefighter's possibility space is an empirical question, not an
algebraic one. What is universal is the structure: that these four
operations are well-defined before truth exists, that they preserve the
pre-logical constraints ($i > 0$, $j < 1$), and that they move the
state toward collapse.
\end{remark}

%======================================================================
\section{Collapse as Algebraic Operation}
\label{sec:collapse}
%======================================================================

\begin{definition}[Collapse]
\label{def:collapse}
The \textbf{collapse operator} $\mathbf{j}: \hat{V} \to V$ is the map
that assigns a truth value to a pre-logical state:
\begin{equation}
\mathbf{j}(?, i, j, k) = (r, \; i', \; 1, \; k')
\end{equation}
where $r \in [0,1]$ is the crystallized truth, $i' \leq i$ (potentiality
is reduced), $j = 1$ (verification is complete), and $k' \leq k$
(selection narrows to the chosen path). We write $\hat{V}_{\mathrm{pre}}$
for the subset of $\hat{V}$ with $r = {?}$.
\end{definition}

The collapse is not instantaneous in general. We distinguish three stages:

\begin{definition}[Stages of collapse]
\label{def:stages}
\leavevmode
\begin{enumerate}[nosep]
\item \textbf{Pre-collapse}: $(?, \; i, \; j_0, \; k_0)$ with $j_0 \approx 0$.
  Possibilities are open, no evidence gathered.
\item \textbf{Partial collapse}: $(\tilde{r}, \; i', \; j', \; k')$ with
  $0 < j' < 1$. Some evidence has been gathered; truth is partially
  crystallized ($\tilde{r}$ is a provisional estimate). The state has
  entered $V$ but verification is incomplete.
\item \textbf{Total collapse}: $(r, \; 0, \; 1, \; 0)$. Truth is fully
  crystallized, no potentiality remains, verification is complete,
  selection is resolved. The state is fully logical.
\end{enumerate}
\end{definition}

\begin{theorem}[Properties of collapse]
\label{thm:collapse-properties}
The collapse operator $\mathbf{j}$ satisfies:
\begin{enumerate}[nosep]
\item \textbf{Irreversibility}: Once $r$ is assigned, it cannot be
  retracted to~$?$. The map $\mathbf{j}$ has no left inverse on its image.
\item \textbf{Monotonicity of verification}: Along any collapse trajectory,
  $j(t)$ is non-decreasing.
\item \textbf{Reduction of potentiality}: Along any collapse trajectory,
  $|i(t)|$ is non-increasing.
\item \textbf{Idempotence on logical states}: If $v \in V$ (already
  logical), then $\mathbf{j}(v) = v$. Collapse of a collapsed state is
  the identity.
\end{enumerate}
\end{theorem}

\begin{proof}
(1)~By construction: $? \notin [0,1]$, so no element of $V$ maps back
to~$\hat{V}_{\mathrm{pre}}$.
(2)~Verification accumulates evidence; the constraint $j' \leq j''$ for
$t' < t''$ follows from the fact that each partial evaluation adds
$\epsilon > 0$ (Proposition~\ref{prop:partial-eval}).
(3)~Each partial evaluation reduces potentiality by $\delta \geq 0$, and
each selection reduces it by $\delta' \geq 0$; the total is non-increasing.
(4)~If $r \in [0,1]$, the state is already in $V$, and the collapse
operator acts as the identity on its codomain by definition.
\end{proof}

\begin{remark}[Collapse is parametrised, not predicted]
\label{rem:collapse-param}
The collapse operator $\mathbf{j}$ does not predict \emph{which}
truth value $r$ crystallises: the attractor $\mu_0$ in
Theorem~\ref{thm:wave-fep-equiv} is a parameter determined by the
domain---what the highway object is, what the judge rules, what the
city builds. The framework characterises the universal
\emph{dynamics} of collapse (damped oscillation toward $\mu_0$)
and the universal \emph{structure} of the pre-logical regime
(Boolean lattice, four operations), while leaving the
domain-specific \emph{content} (which particular $r$ emerges)
to the domain. This separation is deliberate: the substrate
independence of Section~\ref{sec:substrate} depends on the
process being invariant while the outcome varies.
\end{remark}

%======================================================================
\section{The Damped Wave as Free Energy Minimization}
\label{sec:wave-fep}
%======================================================================

The transition from pre-logical to logical is not merely algebraic---it
has dynamics. We prove that the collapse trajectory is governed by a
damped harmonic oscillator that emerges exactly from the free energy
principle.

\begin{theorem}[DHO = FEP]
\label{thm:wave-fep-equiv}
Under the Laplace approximation with generalized coordinates
$\tilde{\mu} = (\mu, \mu')$ and quadratic free energy
\begin{equation}
F(\mu, \mu') = \tfrac{1}{2}\kappa(\mu - \mu_0)^2
  + \tfrac{1}{2}\rho(\mu')^2
\end{equation}
the Friston dynamics~\eqref{eq:friston} with dissipation
$\Gamma = \operatorname{diag}(0, \gamma)$ and solenoidal coupling
$Q = \bigl(\begin{smallmatrix} 0 & -\sigma \\ \sigma & \phantom{-}0
\end{smallmatrix}\bigr)$ yield:
\begin{equation}
\label{eq:dho}
\ddot{\mu} + \gamma\rho\,\dot{\mu}
  + \sigma\kappa(1 + \sigma\rho)\,(\mu - \mu_0) = 0
\end{equation}
whose underdamped solution
($\sigma\kappa(1+\sigma\rho) > (\gamma\rho)^2/4$) is:
\begin{equation}
\label{eq:solution}
\mu(t) = A e^{-dt}\sin(\omega_d\, t + \phi) + \mu_0
\end{equation}
with $d = \gamma\rho/2$ and
$\omega_d = \sqrt{\sigma\kappa(1+\sigma\rho) - (\gamma\rho)^2/4}$.
\end{theorem}

\begin{proof}
The free energy gradient is
$\nabla_{\tilde{\mu}} F = \bigl(\kappa(\mu-\mu_0),\;\rho\mu'\bigr)^\top$.
The shift operator and flow matrices are:
\[
D = \begin{pmatrix} 0 & 1 \\ 0 & 0 \end{pmatrix}, \qquad
\Gamma + Q = \begin{pmatrix} 0 & -\sigma \\ \sigma & \gamma \end{pmatrix}
\]

Applying Eq.~\eqref{eq:friston}:
\begin{align}
\dot{\mu}  &= \mu' - (-\sigma\rho\mu')
            = (1 + \sigma\rho)\,\mu'
            \label{eq:flow1} \\
\dot{\mu'} &= -\sigma\kappa(\mu - \mu_0) - \gamma\rho\,\mu'
            \label{eq:flow2}
\end{align}

From \eqref{eq:flow1}: $\mu' = \dot{\mu}/(1+\sigma\rho)$.
Differentiating: $\dot{\mu'} = \ddot{\mu}/(1+\sigma\rho)$.
Substituting into \eqref{eq:flow2} and multiplying by $(1+\sigma\rho)$:
\begin{equation}
\ddot{\mu} + \gamma\rho\,\dot{\mu}
  + \sigma\kappa(1+\sigma\rho)\,(\mu - \mu_0) = 0
\end{equation}

This is a damped harmonic oscillator with effective damping
$b_{\mathrm{eff}} = \gamma\rho$ and effective spring constant
$k_{\mathrm{eff}} = \sigma\kappa(1+\sigma\rho)$. The characteristic
equation $\lambda^2 + \gamma\rho\,\lambda + \sigma\kappa(1+\sigma\rho) = 0$
has discriminant:
\begin{equation}
\Delta = (\gamma\rho)^2 - 4\sigma\kappa(1+\sigma\rho)
\end{equation}
When $\Delta < 0$, the roots are complex conjugates
$\lambda = -d \pm i\omega_d$, yielding the underdamped
solution~\eqref{eq:solution}.
\end{proof}

\begin{remark}[Necessity of solenoidal coupling]
\label{rem:solenoidal}
With $Q = 0$ (purely dissipative dynamics) and $\Gamma = I$, the same
quadratic~$F$ yields
$\ddot{\mu} + (\kappa + \rho)\dot{\mu} + \rho\kappa(\mu - \mu_0) = 0$,
whose discriminant $(\kappa-\rho)^2 \geq 0$ is always non-negative.
This system is \emph{always} overdamped---it can never oscillate.
Solenoidal coupling ($\sigma > 0$) is therefore necessary for the
underdamped regime within the FEP framework.
Physically: the solenoidal component represents circulation in phase
space---the system does not simply descend the free energy gradient, it
\emph{spirals} around the attractor while descending.
\end{remark}

\begin{remark}[Interpretation for collapse dynamics]
\label{rem:collapse-dynamics}
In the pre-logical context, the FEP parameters receive a natural
interpretation:
\begin{itemize}[nosep]
\item $\mu$: the agent's current estimate (partial collapse state).
\item $\mu_0$: the attractor---the truth that will eventually crystallize.
\item $A = D_{KL}(q_0 \| p)$: initial divergence between belief and reality.
\item $d = \gamma\rho/2$: rate of surprise reduction.
\item $\omega_d$: frequency of oscillation between over- and under-estimates.
\item $\sigma$: strength of the solenoidal coupling that enables oscillation.
\end{itemize}
Collapse is not a sudden jump but a \emph{damped oscillation}:
the agent overshoots, corrects, overshoots less, corrects again, spiraling
toward the truth.
\end{remark}

\subsection{Parameter mapping}

\begin{table}[H]
\centering
\caption{Parameter mapping between the wave model and the FEP.
Each row is a definitional identification under the derivation, not a
loose analogy.}
\label{tab:params}
\begin{tabularx}{\textwidth}{l Y Y}
\toprule
\textbf{Wave} & \textbf{FEP equivalent} & \textbf{Interpretation} \\
\midrule
$A$ (amplitude) & $D_{KL}(q_0 \| p)$ & Initial divergence between
  belief and reality \\
$d$ (damping) & $\gamma\rho/2$ & Dissipative precision on velocity;
  rate of surprise reduction \\
$f$ (frequency) & $\frac{1}{2\pi}\sqrt{\sigma\kappa(1{+}\sigma\rho)
  - (\gamma\rho)^2/4}$ & Solenoidal--curvature interaction \\
$\phi$ (phase) & Prior mean offset & Initial belief about the state \\
$\mu_0$ (zero line) & Attractor & The state the system expects to
  occupy \\
$Ae^{-dt}$ (envelope) & $F(t) \to F_{\min}$ & Free energy approaching
  its minimum \\
$\sigma$ (new) & Solenoidal coupling & Circulation in phase space;
  enables oscillation \\
\bottomrule
\end{tabularx}
\end{table}

\subsection{Conditions for oscillation}

The system oscillates when:
\begin{equation}
\sigma\kappa(1 + \sigma\rho) > \frac{(\gamma\rho)^2}{4}
\end{equation}

\begin{corollary}[Three dynamical regimes]
\label{cor:regimes}
\leavevmode
\begin{enumerate}[nosep]
\item \textbf{Underdamped}
  ($\sigma\kappa(1{+}\sigma\rho) > (\gamma\rho)^2/4$):
  Oscillation. The system spirals around the attractor with decaying
  amplitude.
\item \textbf{Critically damped}
  ($\sigma\kappa(1{+}\sigma\rho) = (\gamma\rho)^2/4$):
  Fastest convergence without oscillation.
\item \textbf{Overdamped}
  ($\sigma\kappa(1{+}\sigma\rho) < (\gamma\rho)^2/4$):
  Monotone exponential decay. No oscillation.
\end{enumerate}
\end{corollary}

%======================================================================
\section{Two Routes to Oscillation}
\label{sec:two-routes}
%======================================================================

\begin{proposition}[Geometric route to oscillation]
\label{prop:geometric}
Under the standard Friston dynamics ($\Gamma = I$, $Q = 0$) with a
\emph{non-diagonal} free energy that includes a cross-term~$\eta$:
\begin{equation}
\label{eq:F-eta}
F_\eta(\mu, \mu') = \tfrac{1}{2}\kappa(\mu - \mu_0)^2
  + \tfrac{1}{2}\rho(\mu')^2 + \eta(\mu - \mu_0)\mu'
\end{equation}
the dynamics yield:
\begin{equation}
\label{eq:dho-eta}
\ddot{\mu} + (\kappa + \rho)\,\dot{\mu}
  + [\rho\kappa + \eta(1-\eta)]\,(\mu - \mu_0) = 0
\end{equation}
with discriminant
$\Delta = (\kappa - \rho)^2 - 4\eta(1-\eta)$.
For $0 < \eta < 1$ and $(\kappa - \rho)^2 < 4\eta(1-\eta)$, the system
is underdamped---oscillation arises \emph{without} solenoidal coupling.
\end{proposition}

\begin{proof}
The gradient is
$\nabla F_\eta = \bigl(\kappa(\mu{-}\mu_0) + \eta\mu',\;
\rho\mu' + \eta(\mu{-}\mu_0)\bigr)^\top$.
Standard Friston dynamics
($\dot{\tilde{\mu}} = D\tilde{\mu} - \nabla F_\eta$):
\begin{align}
\dot{\mu}  &= (1-\eta)\,\mu' - \kappa(\mu - \mu_0) \\
\dot{\mu'} &= -\rho\,\mu' - \eta(\mu - \mu_0)
\end{align}
From the first equation:
$\mu' = \bigl(\dot{\mu} + \kappa(\mu{-}\mu_0)\bigr)/(1-\eta)$.
Differentiating and substituting:
\[
\ddot{\mu} + (\kappa + \rho)\,\dot{\mu}
  + [\rho\kappa + \eta(1{-}\eta)]\,(\mu - \mu_0) = 0
\]
The discriminant is
$(\kappa{+}\rho)^2 - 4[\rho\kappa + \eta(1{-}\eta)]
= (\kappa{-}\rho)^2 - 4\eta(1{-}\eta)$, which is negative when
$(\kappa{-}\rho)^2 < 4\eta(1{-}\eta)$.
\end{proof}

\begin{remark}[Two independent routes]
\label{rem:two-routes}
Theorem~\ref{thm:wave-fep-equiv} and Proposition~\ref{prop:geometric}
provide two independent routes to oscillation within the FEP:
\begin{enumerate}[nosep]
\item \textbf{Dynamical route}: diagonal~$F$ with solenoidal coupling
  ($\sigma > 0$). The system spirals due to non-dissipative flow.
\item \textbf{Geometric route}: non-diagonal~$F$ with standard dynamics
  ($Q = 0$). The free energy landscape itself drives oscillation.
\end{enumerate}
The geometric route is physically suggestive: the cross-term~$\eta$
represents coupling between a system's state and its rate of change
within the free energy functional. In the pre-logical context, this
means that \emph{what the agent currently believes} and \emph{how fast
that belief is changing} are coupled within the free energy landscape
---the shape of uncertainty itself can drive oscillatory collapse.
\end{remark}

%======================================================================
\section{Information as Compound Concept}
\label{sec:information}
%======================================================================

The duality synthesis framework \citep{p12} shows that opposing poles
generate a compound concept through
$a \oplus \bar{a} = \mathrm{lcm}(\Phi(a), \Phi(\bar{a}))$. We apply
this to the fundamental duality of this paper.

\begin{proposition}[Information as compound concept]
\label{thm:information}
The duality (pre-logical, logical) generates \textbf{information} as its
compound concept. Specifically:
\begin{enumerate}[nosep]
\item The two poles are genuinely opposing: the pre-logical state
  $(?, i, j_0, k_0)$ has maximal potentiality and no truth; the logical
  state $(r, 0, 1, 0)$ has crystallised truth and no potentiality.
  Neither is reducible to the other.
\item Information requires \emph{both} poles: it is neither pure
  possibility (that is ignorance) nor pure truth (that is fact), but
  the \emph{reduction} of possibility \emph{toward} truth.
\item The verification axis~$j$ measures the degree of collapse between
  the poles. The selection axis~$k$ encodes the agent's commitment.
\end{enumerate}
\end{proposition}

\begin{proof}
This is a structural argument, not a computation in the prime algebra:
the pre-logical/\allowbreak logical duality is a \emph{meta}-duality (a duality of
the framework itself, not of concepts within the 72-primitive ontology),
so the prime signatures $\Phi(\text{pre-logical})$ and
$\Phi(\text{logical})$ are not defined. What we verify is the
\emph{pattern} of synthesis from~\citet{p12}.

\textbf{Opposition.}
The pre-logical state has $r = {?}$ (undefined), $i > 0$ (open
possibilities), $j < 1$ (incomplete verification). The logical state has
$r \in [0,1]$ (crystallised), $i = 0$ (no remaining possibilities),
$j = 1$ (complete verification). On every axis where one pole is maximal,
the other is minimal. This is the strongest form of opposition in the
quaternionic framework.

\textbf{Irreducibility.}
Pure possibility ($r$ absent) cannot derive truth; pure truth ($i = 0$)
cannot recover the space of alternatives that was collapsed.
The concept ``information'' requires reference to both: information
is the measure of how much possibility was eliminated to reach truth.

\textbf{Compatibility with Shannon.}
Shannon entropy $H = -\sum p_i \log p_i$~\citep{shannon1948} measures
information within the $r$~axis: probabilities are truth values, and
$I = H_{\text{before}} - H_{\text{after}}$ is the reduction from
uncertainty to certainty. This is the restriction of the pre-logical/\allowbreak logical
transition to the regime where $r$ is defined (indeterminacy, not
indefiniteness). The quaternionic framework does not derive Shannon
entropy; it extends the concept to the regime where $r$ is not yet
defined, with geodesic distance on~$S^3$ replacing entropy as the
measure of separation between current state and attractor.
\end{proof}

\begin{remark}[The bit as elementary collapse]
A single bit of information corresponds to one step of collapse: the
resolution of one binary question (``is the object a tree? yes/no'').
In the prime lattice, this is a Hamming-distance-1 move. In the FEP
dynamics, it is one half-period of the damped oscillation (overshoot,
then correction). The bit is the quantum of the pre-logical-to-logical
transition.
\end{remark}

%======================================================================
\section{Substrate Independence}
\label{sec:substrate}
%======================================================================

The structure (pre-logical state $\to$ collapse $\to$ logical state)
is substrate-independent: the same algebraic framework applies across
domains that differ in their physical instantiation. We demonstrate this
with four examples.

\subsection{Quantum measurement}

\begin{domainbox}{blue}
\textbf{Domain: Quantum physics --- Schr\"odinger's cat.}

\begin{tabularx}{\textwidth}{l Y}
\toprule
\textbf{Component} & \textbf{Instantiation} \\
\midrule
Pre-logical state & Closed box; Re undefined \\
Potentiality ($i$) & Alive and dead as superposed possibilities \\
Verification ($j$) & Opening the box / measuring \\
Selection ($k$) & Design of the experiment \\
Collapse & Measurement $\to$ eigenvalue \\
Logical state & Alive or dead---crystallized by observation \\
\bottomrule
\end{tabularx}

The quantum case is often presented as exotic. It is not. The firefighter
facing an unidentified highway obstruction is in the same algebraic
situation: a superposition of possibilities that collapses upon
observation. The difference is substrate (quantum amplitudes vs.\
macroscopic objects), not structure.
\end{domainbox}

\subsection{Legal reasoning}

\begin{domainbox}{red}
\textbf{Domain: Law --- Case without precedent.}

\begin{tabularx}{\textwidth}{l Y}
\toprule
\textbf{Component} & \textbf{Instantiation} \\
\midrule
Pre-logical state & Novel case; no applicable jurisprudence \\
Potentiality ($i$) & Possible legal arguments \\
Verification ($j$) & Judicial ruling \\
Selection ($k$) & Lawyer's litigation strategy \\
Collapse & Ruling $\to$ new precedent \\
Logical state & Jurisprudence created---crystallized by the court \\
\bottomrule
\end{tabularx}

The legal method IRAC (Issue, Rule, Application, Conclusion) maps
directly onto the pre-logical framework:
\emph{Issue}~$=$ identify which Re does not yet exist (the pre-logical state);
\emph{Rule}~$=$ existing Re (crystallized law);
\emph{Application}~$=$ $j + k$ (verify applicability, select strategy);
\emph{Conclusion}~$=$ collapse $\to$ new Re.

Each law can be viewed as a quaternion: $r$ is the literal text, $i$ the
possible interpretations, $j$ the case law that verifies it, $k$ the
selection of which precedent applies. Changing jurisdiction (e.g.,
from one state to another) is a rotation $q$ on $S^3$: the commutation
theorem predicts that laws depending only on Re (constitutional
provisions) are invariant under rotation, while laws depending on Im
(local interpretation) vary.
\end{domainbox}

\subsection{Emergency response}

\begin{domainbox}{orange}
\textbf{Domain: Emergencies --- Unidentified highway obstruction.}

\begin{tabularx}{\textwidth}{l Y}
\toprule
\textbf{Component} & \textbf{Instantiation} \\
\midrule
Pre-logical state & Vague report: ``large object on Highway~45'' \\
Potentiality ($i$) & Tree, vehicle, animal, boulder \\
Verification ($j$) & Arriving at the scene \\
Selection ($k$) & Equipment choice \\
Collapse & Visual identification $\to$ confirmed object \\
Logical state & ``It is a fallen tree''---crystallized by observation \\
\bottomrule
\end{tabularx}

The optimal decision under a pre-logical state is the minimax preparation
(Proposition~\ref{prop:minimax}): carry equipment that minimizes the
maximum distance to all possible outcomes. The firefighter brings crane,
hammers, \emph{and} winch because, regardless of what the object is, the
preparation should be close to the solution.

Partial evaluation is available: send a drone first ($j$ partial, low
cost) before committing the full crew ($j$ total, high cost). The
sequence of partial collapses is an optimization over the lattice of
possibilities.
\end{domainbox}

\subsection{Urban planning}

\begin{domainbox}{green!50!black}
\textbf{Domain: Urbanism --- City 2040.}

\begin{tabularx}{\textwidth}{l Y}
\toprule
\textbf{Component} & \textbf{Instantiation} \\
\midrule
Pre-logical state & Uncertain future: water, transit, growth \\
Potentiality ($i$) & Scenarios for population, climate, economy \\
Verification ($j$) & Construction (irreversible) \\
Selection ($k$) & Budget allocation \\
Collapse & Building $\to$ physical infrastructure \\
Logical state & City built---crystallized by construction \\
\bottomrule
\end{tabularx}

Each urban decision collapses possibilities: building a highway
\emph{here} eliminates a park \emph{there}. Each $j$ that is executed
closes branches of the adjacent possible. The lattice formalism
(Section~\ref{sec:lattice}) makes this precise: the known set
$K$ grows, but the adjacent possible $AP(K)$ changes structure with
each commitment.

The irreversibility of construction makes urban planning a
particularly pure example of collapse: unlike quantum measurement
(which can be repeated) or legal rulings (which can be overturned on
appeal), a building \emph{cannot be unbuilt} at reasonable cost.
\end{domainbox}

\subsection{Structural invariance}

\begin{table}[H]
\centering
\caption{Substrate independence: the same structure across four domains.}
\label{tab:substrate}
\begin{tabularx}{\textwidth}{l Y Y Y Y}
\toprule
& \textbf{Quantum} & \textbf{Legal} & \textbf{Emergency} & \textbf{Urban} \\
\midrule
$r$ undefined & Superposition & No precedent & Unknown object & Unknown future \\
$i > 0$ & Amplitudes & Arguments & Possibilities & Scenarios \\
$j$ (collapse) & Measurement & Ruling & Observation & Construction \\
$k$ (selection) & Experiment & Strategy & Equipment & Budget \\
Lattice & Hilbert space & Case space & Object space & Scenario space \\
Irreversibility & Wavefunction & Partial & Low & Total \\
\bottomrule
\end{tabularx}
\end{table}

The structural invariance across these four domains---differing in
substrate, time scale, reversibility, and social context---supports the
claim that pre-logical states are not a special case but the
\emph{generic} situation for reasoning agents. Logic (the assignment and
manipulation of truth values) is the endpoint of a process that begins
in the pre-logical regime.

%======================================================================
\section{Discussion}
\label{sec:discussion}
%======================================================================

\subsection{Step zero}

Classical epistemology describes a chain:
(1)~opposites exist (hot/cold),
(2)~a scale connects them,
(3)~the zero point is arbitrary (Kelvin, Fahrenheit, Celsius are
rotations on $S^3$),
(4)~a compound concept emerges (temperature),
(5)~perspective and context rotate the reading.

The pre-logical state is \textbf{step zero}: before the agent knows
\emph{which} opposites apply. The firefighter does not know whether the
problem is heavy/light, organic/inorganic, mobile/fixed. The lawyer
does not know whether the case is civil/criminal, state/federal. The
urban planner does not know whether the problem is water, drainage, or
transit.

Verification ($j$) moves the agent from step zero to step one:
collapsing not only truth but \emph{which framework of opposites
applies}. The pre-logical state is the pre-structure of structure---navigating
the adjacent possible of dualities, not of concepts.

\subsection{Relation to probability theory}

Pre-logical states are not probabilistic states. Probability theory
assigns a distribution $P(r)$ over truth values; it handles
\emph{indeterminacy} (the truth exists but is unknown). The pre-logical
state handles \emph{indefiniteness} (no truth has crystallized). The
distinction matters:
\begin{itemize}[nosep]
\item Under indeterminacy, Bayesian updating is optimal.
\item Under indefiniteness, the optimal strategy is minimax
  (Proposition~\ref{prop:minimax})---minimize worst-case distance, not
  expected distance.
\end{itemize}
The two regimes connect through the FEP: Bayesian updating corresponds
to the underdamped and overdamped regimes (convergence to an existing
attractor), while transitions between attractors correspond to shifts
in the pre-logical landscape---new primes entering the lattice.

\subsection{Limitations}

\begin{enumerate}[nosep]
\item \textbf{Structural, not quantitative.} The substrate independence
  demonstrations establish that the same algebraic framework applies
  across four domains, but they do not derive numerical predictions for
  any specific domain. The value is in the structural correspondence:
  identifying which operations are available and which properties hold,
  not in predicting particular outcomes.
\item \textbf{Collapse endpoint is domain-specific.} The FEP dynamics
  describe \emph{how} collapse occurs (damped oscillation toward the
  attractor~$\mu_0$) but the attractor itself is a parameter determined
  by the domain---what the object on the highway actually is, what the
  judge actually rules. This is by design: the framework characterises
  the universal \emph{process} of collapse while leaving the
  domain-specific \emph{content} to the domain.
\item \textbf{Four domains is few.} The substrate independence claim
  would be strengthened by additional domains (medicine, AI/ML,
  financial markets).
\item \textbf{Conservative lattice.} The prime product lattice allows
  all combinations; real domains have constraints. The lattice is the
  ``ideal gas'' of the adjacent possible---exact and tractable, but
  missing thermodynamic constraints.
\end{enumerate}

%======================================================================
\section{Conclusion}
\label{sec:conclusion}
%======================================================================

Logic is not the beginning of reasoning---it is the end.

Before truth exists, there is navigation: possibilities that expand
(Theorem~\ref{thm:expansion}), paths that depend on history
(Theorem~\ref{thm:path}), exploration that falls exponentially behind
the growth of the possible (Theorem~\ref{thm:faster}), and futures
that cannot be prestated (Theorem~\ref{thm:nonprestate}). The agent
gathers evidence (partial evaluation), prepares for all outcomes
(minimax preparation), and commits (selection)---all without a truth value
to operate on.

The transition from this pre-logical regime to the logical one---collapse---is
a damped oscillation governed by the free energy principle
(Theorem~\ref{thm:wave-fep-equiv}), with solenoidal coupling necessary
for the system to spiral rather than slide to its attractor
(Remark~\ref{rem:solenoidal}). Two independent routes to oscillation
exist: dynamical (solenoidal flow) and geometric (cross-coupled free
energy) (Proposition~\ref{prop:geometric}).

The duality between pre-logical and logical generates information as its
compound concept (Proposition~\ref{thm:information}): not possibility
alone (that is ignorance) and not truth alone (that is fact), but the
\emph{reduction} of the one toward the other. Shannon entropy operates
within the truth axis (where $r$ is defined); the quaternionic framework
extends the concept to the regime where $r$ is not yet defined, with
geodesic distance replacing entropy as the measure of separation.

The structure is substrate-independent: quantum measurement, legal
reasoning, emergency response, and urban planning all instantiate the
same algebraic framework with different physical substrates. This
suggests that the pre-logical state is not an edge case but the generic
condition of reasoning agents---and that classical logic describes the
world after the measurement, not before it.

%======================================================================
% BIBLIOGRAPHY
%======================================================================

\bibliographystyle{plainnat}

\begin{thebibliography}{99}

\bibitem[Buckley et al.(2017)]{buckley2017}
Buckley, C.L., Kim, C.S., McGregor, S., \& Seth, A.K. (2017).
The free energy principle for action and perception: A mathematical review.
\textit{Journal of Mathematical Psychology}, 81, 55--79.

\bibitem[Friston(2010)]{friston2010}
Friston, K. (2010).
The free-energy principle: a unified brain theory?
\textit{Nature Reviews Neuroscience}, 11(2), 127--138.

\bibitem[Kauffman(2000)]{kauffman2000}
Kauffman, S.A. (2000).
\textit{Investigations}.
Oxford University Press.

\bibitem[Loreto et al.(2017)]{loreto2017}
Loreto, V., Servedio, V.D.P., Strogatz, S.H., \& Tria, F. (2017).
Dynamics on expanding spaces: modeling the emergence of novelties.
In M.\ Degli Esposti, E.G.\ Altmann, \& F.\ Pachet (Eds.),
\textit{Creativity and Universality in Language}, pp.\ 59--83.
Springer.

\bibitem[Ornelas Brand(2026a)]{p1}
Ornelas Brand, J.A. (2026).
Triadic Neurosymbolic Engine: Prime Factorization as a Neurosymbolic
Bridge: Projecting Continuous Embeddings into Discrete Algebraic Space
for Deterministic Verification.
\textit{Zenodo}.
\href{https://doi.org/10.5281/zenodo.19205804}{doi:10.5281/zenodo.19205804}

\bibitem[Ornelas Brand(2026b)]{p4}
Ornelas Brand, J.A. (2026).
Toward a Layered Ontology of Semantic Duality: 14 Candidate Dualities
Across 6 Algebraic Layers with Preliminary Domain and Neural Evidence
(v1.0).
\textit{Zenodo}.
\href{https://doi.org/10.5281/zenodo.19375166}{doi:10.5281/zenodo.19375166}

\bibitem[Ornelas Brand(2026c)]{p11}
Ornelas Brand, J.A. (2026).
Quaternionic Logic: A G-Lattice Unifying Boolean and Fuzzy Frameworks (v0.1.0).
\textit{Zenodo}.
\href{https://doi.org/10.5281/zenodo.19560986}{doi:10.5281/zenodo.19560986}

\bibitem[Ornelas Brand(2026d)]{p12}
Ornelas Brand, J.A. (2026).
Duality Synthesis in Quaternionic Logic: How Opposites Generate Truth (v0.1.0).
\textit{Zenodo}.
\href{https://doi.org/10.5281/zenodo.19561633}{doi:10.5281/zenodo.19561633}

\bibitem[Shannon(1948)]{shannon1948}
Shannon, C.E. (1948).
A mathematical theory of communication.
\textit{Bell System Technical Journal}, 27(3), 379--423.

\end{thebibliography}

\end{document}
